\documentclass[10pt,a4paper]{amsart}
\usepackage{amsmath,amssymb,amsthm,graphicx}
\usepackage{verbatim}
\newenvironment{reference}{\comment}{\endcomment}
%\newenvironment{reference}{}{}
\newenvironment{slogan}{\comment}{\endcomment}
\newenvironment{history}{\comment}{\endcomment}

% For commutative diagrams we use Xy-pic
\usepackage[all]{xy}

% We use 2cell for 2-commutative diagrams.
\xyoption{2cell}
% Theorem environments.
%
\theoremstyle{plain}
\newtheorem{theorem}[subsection]{Theorem}
\newtheorem{proposition}[subsection]{Proposition}
\newtheorem{lemma}[subsection]{Lemma}

\theoremstyle{definition}
\newtheorem{definition}[subsection]{Definition}
\newtheorem{example}[subsection]{Example}
\newtheorem{exercise}[subsection]{Exercise}
\newtheorem{situation}[subsection]{Situation}

\theoremstyle{remark}
\newtheorem{remark}[subsection]{Remark}
\newtheorem{remarks}[subsection]{Remarks}

\numberwithin{equation}{subsection}

% Macros
%
\def\lim{\mathop{\mathrm{lim}}\nolimits}
\def\colim{\mathop{\mathrm{colim}}\nolimits}
\def\Spec{\mathop{\mathrm{Spec}}}
\def\Hom{\mathop{\mathrm{Hom}}\nolimits}
\def\Ext{\mathop{\mathrm{Ext}}\nolimits}
\def\SheafHom{\mathop{\mathcal{H}\!\mathit{om}}\nolimits}
\def\SheafExt{\mathop{\mathcal{E}\!\mathit{xt}}\nolimits}
\def\Sch{\mathit{Sch}}
\def\Mor{\mathop{\mathrm{Mor}}\nolimits}
\def\Ob{\mathop{\mathrm{Ob}}\nolimits}
\def\Sh{\mathop{\mathit{Sh}}\nolimits}
\def\NL{\mathop{N\!L}\nolimits}
\def\CH{\mathop{\mathrm{CH}}\nolimits}
\def\proetale{{pro\text{-}\acute{e}tale}}
\def\etale{{\acute{e}tale}}
\def\QCoh{\mathit{QCoh}}
\def\Ker{\mathop{\mathrm{Ker}}}
\def\Im{\mathop{\mathrm{Im}}}
\def\Coker{\mathop{\mathrm{Coker}}}
\def\Coim{\mathop{\mathrm{Coim}}}

% Boxtimes
%
\DeclareMathSymbol{\boxtimes}{\mathbin}{AMSa}{"02}

%
% Macros for moduli stacks/spaces
%
\def\QCohstack{\mathcal{QC}\!\mathit{oh}}
\def\Cohstack{\mathcal{C}\!\mathit{oh}}
\def\Spacesstack{\mathcal{S}\!\mathit{paces}}
\def\Quotfunctor{\mathrm{Quot}}
\def\Hilbfunctor{\mathrm{Hilb}}
\def\Curvesstack{\mathcal{C}\!\mathit{urves}}
\def\Polarizedstack{\mathcal{P}\!\mathit{olarized}}
\def\Complexesstack{\mathcal{C}\!\mathit{omplexes}}
% \Pic is the operator that assigns to X its picard group, usage \Pic(X)
% \Picardstack_{X/B} denotes the Picard stack of X over B
% \Picardfunctor_{X/B} denotes the Picard functor of X over B
\def\Pic{\mathop{\mathrm{Pic}}\nolimits}
\def\Picardstack{\mathcal{P}\!\mathit{ic}}
\def\Picardfunctor{\mathrm{Pic}}
\def\Deformationcategory{\mathcal{D}\!\mathit{ef}}
\usepackage{cleveref}
\crefname{originallemma}{original lemma}{original lemmas}
\Crefname{originallemma}{Original lemma}{Original lemmas}
\makeatletter
\newlabel{lemma-simple-module-unique}{{074E}{source}{}{sourcepointer-1}{}}
\newlabel{lemma-simple-module-unique@cref}{{[originallemma][0][]074E}{[1][0][]source}{}{}{}}
\makeatother
\crefname{originallemma}{original lemma}{original lemmas}
\Crefname{originallemma}{Original lemma}{Original lemmas}
\makeatletter
\newlabel{lemma-tensor-simple}{{074F}{source}{}{sourcepointer-2}{}}
\newlabel{lemma-tensor-simple@cref}{{[originallemma][0][]074F}{[1][0][]source}{}{}{}}
\makeatother
\title{The centralizer of a simple subalgebra of a finite central simple algebra}
\author{The Stacks Project Authors}
\date{stacks-project:brauer.tex:local-snapshot-20261012; GNU FDL 1.2 or later}
\begin{document}
\maketitle
\noindent\textit{Editorial note.} All delivered mathematical text is the author's own, copied from the pinned Stacks Project chapter file sites-cohomology.tex; no statement or proof step was rewritten, reordered or removed.\par\medskip
\noindent\textit{Editorial note.} The proof is the author's complete argument; it states at the outset that it uses the results of a named lemma freely, and every identification and dimension formula it needs is displayed or attributed inside the counted span.\par\medskip
\noindent\textit{Editorial note.} Every labelled reference occurring inside the counted spans is carried in assembly.external\_references as a non-load-bearing navigation pointer (same-chapter ones at their definition site, cross-chapter ones at the invocation inside the counted span), so no load-bearing mathematics is replaced by a stub. No bibliography citation occurs inside the counted spans.\par\medskip
\noindent\textit{Editorial note.} The same-name check for this card scanned every staged Stacks chapter of this repository -- all sources/stacks-*/pinned-original-source.tex files, the card's own chapter excluded -- and the labels found are recorded in dedup.same\_name\_cross\_chapter on the card.\par\medskip
\noindent\textit{Editorial note.} The preamble is the chapter's own retained minimal preamble (sources/stacks-sites-cohomology/retained-preamble.tex), a verbatim excerpt of the native preamble: the Xy-pic setup and the native theorem declarations with the mathematical macros.\par\medskip
\noindent\textit{Allowed standard prerequisites (editorial).}\begin{itemize}
\item The simple module and its endomorphism skew field. Same chapter, Lemma lemma-simple-module-unique. Supplies, for a chosen simple module, that its endomorphism ring is a skew field with the same center acting on the module, that the algebra is the endomorphism ring of the module over that skew field, and the dimension formula used at the end.
\item Simplicity of a tensor product. Same chapter, Lemma lemma-tensor-simple. Supplies that the tensor product of the subalgebra with the opposite of the endomorphism skew field is simple, which makes the centralizer simple by the same lemma on simple modules.
\item Matrix algebras over a skew field. Same chapter, Lemma lemma-matrix-algebras and Lemma lemma-simple-module-unique (4). Supply the description of the tensor product as a matrix algebra over a skew field, and of the module as a direct sum of copies of the simple module of that matrix algebra.
\end{itemize}\medskip
\ifdefined\hypertarget\hypertarget{sourcepointer-1}{}\fi
\noindent\textit{Source reference.} Original lemma 074E: sources/stacks-brauer/pinned-original-source.tex, source lines 498--498. Navigation pointer invoked at proof line 711; the counted span contains the pointer only and the cited result is not reproduced inside it, so the pointer is registered as a non-load-bearing navigation stub in the established form for Stacks material. No bibliography citation is involved.\par\medskip
\ifdefined\hypertarget\hypertarget{sourcepointer-2}{}\fi
\noindent\textit{Source reference.} Original lemma 074F: sources/stacks-brauer/pinned-original-source.tex, source lines 538--538. Navigation pointer invoked at proof line 715; the counted span contains the pointer only and the cited result is not reproduced inside it, so the pointer is registered as a non-load-bearing navigation stub in the established form for Stacks material. No bibliography citation is involved.\par\medskip
\begin{theorem}
\label{theorem-centralizer}
Let $A$ be a finite central simple algebra over $k$, and let
$B$ be a simple subalgebra of $A$. Then
\begin{enumerate}
\item the centralizer $C$ of $B$ in $A$ is simple,
\item $[A : k] = [B : k][C : k]$, and
\item the centralizer of $C$ in $A$ is $B$.
\end{enumerate}
\end{theorem}

\begin{proof}
Throughout this proof we use the results of
Lemma \ref{lemma-simple-module-unique} freely.
Choose a simple $A$-module $M$. Set $L = \text{End}_A(M)$.
Then $L$ is a skew field with center $k$ which acts on the left on $M$
and $A = \text{End}_L(M)$.
Then $M$ is a right $B \otimes_k L^{op}$-module and
$C = \text{End}_{B \otimes_k L^{op}}(M)$.
Since the algebra $B \otimes_k L^{op}$ is simple by
Lemma \ref{lemma-tensor-simple} we see that $C$ is simple (by
Lemma \ref{lemma-simple-module-unique} again).

\medskip\noindent
Write $B \otimes_k L^{op} = \text{Mat}(m \times m, K)$ for some
skew field $K$ finite over $k$. Then $C = \text{Mat}(n \times n, K^{op})$
if $M$ is isomorphic to a direct sum of $n$ copies of the simple
$B \otimes_k L^{op}$-module $K^{\oplus m}$ (the lemma again). Thus we have
$\dim_k(M) = nm [K : k]$, $[B : k] [L : k] = m^2 [K : k]$,
$[C : k] = n^2 [K : k]$, and $[A : k] [L : k] = \dim_k(M)^2$ (by
the lemma again). We conclude that (2) holds.

\medskip\noindent
Part (3) follows because of (2) applied to $C \subset A$ shows
that $[B : k] = [C' : k]$ where $C'$ is the centralizer of $C$ in $A$
(and the obvious fact that $B \subset C')$.
\end{proof}

\end{document}
